\subsection*{把系数看成两个可以调节的旋钮}
给定两种移动指令 $u=(1,1)$ 与 $v=(1,-1)$，可以分别执行任意实数倍。先固定第二个旋钮，只改变第一个，终点沿直线移动；再改变第二个旋钮，这条直线整体平移。因为两方向不平行，这些直线扫过整个平面。“所有能到达的位置”就是张成空间。

为什么允许负数和非整数系数？若只准走整数步，得到的是格点；若只准非负步，通常得到一个扇形区域。张成允许任意实系数，正是为了让加法与数乘自由进行。定义中的量词决定了我们讨论哪个集合。

\subsection*{能到达与能唯一到达}
要到达 $(4,2)$，横向必须满足 $a+b=4$，纵向必须满足 $a-b=2$。相加相减得到 $a=3,b=1$，也就是目标在新方向下的坐标。若另一对系数到达同一位置，两份表达相减会得到一个非平凡零组合；但两个不平行方向不能这样抵消，所以坐标唯一。

把 $v$ 换成 $2u$。两个旋钮看似独立，实际上只改变同一条线上的位置；而 $u=1u+0v=-1u+1v$ 又说明描述出现冗余。基把两方面合在一起：既张成目标空间，又不给同一向量留下多套坐标。

看网格实验时，先预测两根基箭头靠近平行会怎样，再观察结果。网格塌缩将在行列式章节变成面积归零；坐标不再唯一会在零空间章节重新出现。它们不是互不相关的定义，而是同一现象的不同描述。
